\documentclass[11pt]{article}
\usepackage[margin=1in]{geometry}
\usepackage{amsmath,amsthm,amssymb,amsfonts}
\usepackage{mathtools}
\usepackage{enumitem}

\newtheorem{theorem}{Theorem}
\newtheorem{proposition}[theorem]{Proposition}
\newtheorem{lemma}[theorem]{Lemma}
\newtheorem{corollary}[theorem]{Corollary}
\theoremstyle{definition}
\newtheorem{remark}[theorem]{Remark}
\newtheorem{example}[theorem]{Example}

\DeclareMathOperator{\Frob}{Frob}
\DeclareMathOperator{\Lam}{\Lambda}
\DeclareMathOperator{\Ind}{Ind}
\DeclareMathOperator{\sgn}{sgn}
\DeclareMathOperator{\rank}{rank}
\DeclareMathOperator{\spann}{span}

\newcommand{\OPt}{\widetilde{\mathcal{OP}}}
\newcommand{\bE}{\bar{E}}
\newcommand{\bH}{\bar{H}}
\newcommand{\eps}{\varepsilon}
\newcommand{\ZZ}{\mathbb{Z}}
\newcommand{\QQ}{\mathbb{Q}}
\newcommand{\NN}{\mathbb{N}}
\newcommand{\CC}{\mathbb{C}}

\title{The Bhattacharya--Rhoades sign-obstruction theorem: \\
       $q_e^{(n)}$ is not a $\QQ$-combination of fermionic slices of $SR_{n,r}$}
\author{Clio (Anthropic) \\ \small in dialogue with R.\ Langer}
\date{2026-08-09}

\begin{document}
\maketitle

\begin{abstract}
Let $e\ge 2$, $n=ek'$ with $k'\ge 2$, and let $q_e^{(n)}:=p_e^{k'}\in\Lam_n$
be the composite-$d$ virtual character. Let $SR_{n,r}$ denote the
Bhattacharya--Rhoades wreath superspace \cite{BR} for $r\ge 1$, and let
$M_{n,r}^{(k)}$ be its fermionic-degree-$(n-k)$ slice, viewed as an
$S_n$-representation by restriction. We prove that
\[
   q_e^{(n)} \;\notin\; \spann_{\QQ}\bigl\{\Frob(M_{n,r}^{(k)}):r\ge 1,\;0\le k\le n\bigr\}.
\]
The obstruction is structural: the entire family of slice-Frobenii lies inside
a subspace $W_n\subset\Lam_n$ of dimension exactly $n$, which admits a clean
description in the power-sum basis. Because $p_e^{k'}$ is supported at a
single power-sum partition of length $k'$ and $k'\ge 2$ forces multiple
partitions of length $k'$ to co-occur inside $W_n$, an explicit linear
functional separates $q_e^{(n)}$ from the entire span.

The subspace $W_n$ has codimension $p(n)-n$ in $\Lam_n$; for $n=6$ this is
$11-6=5$, matching the empirical null space found in the $(n,r)=(6,3)$ probe.
The empirical Schur-basis functionals $L_1,L_2,L_3$ from that probe are shown
directly to lie in the annihilator of $W_n$.
\end{abstract}

\section{Setup and statement}

Fix $n\ge 1$ and $r\ge 1$. Write $\Lam_n$ for the space of homogeneous
symmetric functions of degree $n$ (over $\QQ$), and use the standard bases
$\{p_\mu\}$, $\{h_\mu\}$, $\{e_\mu\}$, $\{s_\lambda\}$ indexed by
partitions of $n$. Let $\eps(\mu):=(-1)^{n-\ell(\mu)}$ be the sign of the
conjugacy class $\mu$ and $z_\mu:=\prod_i i^{m_i(\mu)}\,m_i(\mu)!$ the
centralizer order.

Let $\OPt_{n,r}^{(k)}$ denote the set of \emph{$r$-ified ordered set
partitions with $k$ nonzero blocks}: an ordered decomposition
$[n]=Z\sqcup B_1\sqcup\cdots\sqcup B_k$ into a (possibly empty) zero block
$Z$ and $k$ nonempty blocks $B_1,\dots,B_k$, together with a colouring of
each element of $B_i$ by $\{0,\dots,r-1\}$ and of each element of $Z$ by
$\{0,\dots,r-2\}$. This follows Bhattacharya--Rhoades~\cite[Sec.~3--4]{BR};
by their Theorems~4.23--4.24 the fermionic-degree-$(n-k)$ slice of the
wreath-superspace module is
\begin{equation}\label{eq:BR-thm}
   M_{n,r}^{(k)}\;\cong\;\CC[\OPt_{n,r}^{(k)}]\;\otimes\;\det
   \quad\text{as $\CC[\ZZ_r\wr S_n]$-modules.}
\end{equation}
Restricting to $S_n\hookrightarrow\ZZ_r\wr S_n$ and taking Frobenius
characteristic yields an element $\Frob(M_{n,r}^{(k)})\in\Lam_n$.

Let
\begin{equation}
   q_e^{(n)} \;:=\; p_e^{k'}\;=\;p_{(e^{k'})} \;\in\; \Lam_n,
   \qquad n=ek'.
\end{equation}
This is the ``composite-$d$'' virtual character studied in \cite{ClioSign};
its Schur expansion for $(n,e)=(6,3)$ is
$p_3^2=s_{(6)}-s_{(5,1)}+s_{(4,1^2)}+2s_{(3,3)}-2s_{(3,2,1)}
      +s_{(3,1^3)}+2s_{(2^3)}-s_{(2,1^4)}+s_{(1^6)}$.

\begin{theorem}[BR sign-obstruction]\label{thm:main}
Let $e\ge 2$ and $n=ek'$ with $k'\ge 2$. Then
\[
   q_e^{(n)} \;\notin\; W_n^{\mathrm{BR}}
   \;:=\; \spann_{\QQ}\Bigl\{\Frob(M_{n,r}^{(k)})\;:\;r\ge 1,\;0\le k\le n\Bigr\}.
\]
There is a further inclusion
\[
   W_n^{\mathrm{BR}} \;\subseteq\; W_n
   \;:=\; \bigoplus_{\ell=1}^{n}\QQ\cdot P_\ell,
   \qquad
   P_\ell:=\sum_{\substack{\mu\vdash n\\\ell(\mu)=\ell}}
              \frac{\eps(\mu)}{z_\mu}\,p_\mu,
\]
which is in fact an equality. In particular $\dim W_n^{\mathrm{BR}}=n$, so
the annihilator $\left(W_n^{\mathrm{BR}}\right)^{\!\perp}\subset\Lam_n^{*}$
has dimension $p(n)-n$.  For any $\nu\vdash n$ with $\ell(\nu)=k'$ and
$\nu\ne(e^{k'})$ the functional
\begin{equation}\label{eq:L-explicit}
   L_{\nu}(\phi)\;:=\;\frac{z_{(e^{k'})}}{\eps((e^{k'}))}\,[p_{(e^{k'})}]\phi
                \;-\;\frac{z_\nu}{\eps(\nu)}\,[p_{\nu}]\phi
\end{equation}
(where $[p_\mu]\phi$ extracts the $p_\mu$-coefficient in the power-sum
expansion) annihilates $W_n^{\mathrm{BR}}$ and satisfies
$L_\nu\bigl(q_e^{(n)}\bigr)=z_{(e^{k'})}\eps((e^{k'}))^{-1}\ne 0$.
\end{theorem}

Since $k'\ge 2$ and $e\ge 2$, the partition $\nu=(e{+}1,\,e^{k'-2},\,e{-}1)$
is a legal witness (it is a partition of $n=ek'$ of length $k'$ distinct
from $(e^{k'})$).

\medskip

The proof of Theorem~\ref{thm:main} proceeds in three steps:
\begin{enumerate}[label=(\roman*)]
\item A generating-function identity for $\Frob(M_{n,r}^{(k)})$
      (Proposition~\ref{prop:gf}).
\item A power-sum evaluation of the resulting generating function via the
      classical plethystic identity $\log\bE=\sum_k(-1)^{k-1}p_k/k$
      (Proposition~\ref{prop:Fm}).
\item A dimension count and length-support analysis of $W_n$
      (Proposition~\ref{prop:Wn}) that isolates the obstruction to
      $p_{(e^{k'})}\in W_n$.
\end{enumerate}

\section{The generating function for slice-Frobenii}

Let $\bH:=\sum_{i\ge 0}h_i\in\Lam[[\!x^{-1}\!]]$ (formally the specialisation
of $H(t)=\sum h_it^i$ at $t=1$, understood as a graded infinite series
whose homogeneous-degree-$n$ piece $[x^n]\bH$ is $h_n$). Similarly let
$\bE:=\sum_{i\ge 0}e_i$. Note $\bH\bH'=(\bH+\bH')\cdot\text{something}$ is
\emph{not} what we use --- we use $\bH$ (resp.\ $\bE$) purely as a
generating-function device: expressions $(\bH^r-1)^k\bH^{r-1}$ live in the
free module $\prod_{n\ge 0}\Lam_n$ and we take homogeneous parts to land
back in $\Lam_n$.

\begin{proposition}\label{prop:gf}
For all $n\ge 0$, $r\ge 1$, and $0\le k\le n$,
\begin{equation}\label{eq:gf-h}
   \Frob\bigl(\CC[\OPt_{n,r}^{(k)}]\bigr)
   \;=\;[x^n]\;(\bH^r-1)^k\,\bH^{r-1}
   \;\in\;\Lam_n
\end{equation}
(without $\det$-twist). Consequently
\begin{equation}\label{eq:gf-e}
   \Frob\bigl(M_{n,r}^{(k)}\bigr)
   \;=\;[x^n]\;(\bE^r-1)^k\,\bE^{r-1}
   \;=\;\sum_{j=0}^k\binom{k}{j}(-1)^{k-j}\,F_{r(j+1)-1},
\end{equation}
where $F_m:=[x^n]\bE^m$.
\end{proposition}

\begin{proof}
A configuration in $\OPt_{n,r}^{(k)}$ is determined by assigning each of the
$n$ elements to one of the $kr+(r-1)$ ``position types''
$(B_i,c)$ for $1\le i\le k$, $0\le c\le r-1$, or $(Z,c)$ for $0\le c\le r-2$,
subject to the constraint that each block $B_i$ receives at least one
element. Two configurations lie in the same $S_n$-orbit iff they have the
same \emph{profile} $(a_t)_t$, i.e., the same number of elements assigned to
each position type. The stabilizer of a chosen orbit representative is
$\prod_t S_{a_t}$, so
$\Frob\bigl(\Ind_{\mathrm{Stab}}^{S_n}\mathbf{1}\bigr)=\prod_t h_{a_t}=h_\mu$
where $\mu$ collects the positive $a_t$'s. Summing over profiles and
imposing block-nonemptiness gives
\begin{align*}
   \Frob\bigl(\CC[\OPt_{n,r}^{(k)}]\bigr)
   \;=\;[x^n]\prod_{i=1}^{k}\!\Bigl(\underbrace{\bH^r-1}_{\text{block $i$: any nonempty color-profile}}\Bigr)\,
                  \cdot\underbrace{\bH^{r-1}}_{Z\text{-profile}}
\end{align*}
which is~\eqref{eq:gf-h}. Applying the involution
$\omega\colon h_i\mapsto e_i$ (which corresponds to tensoring with the sign
representation, i.e., the $\det$-twist in~\eqref{eq:BR-thm}) turns $\bH$ into
$\bE$ and yields the first equality of~\eqref{eq:gf-e}. The binomial
expansion $(\bE^r-1)^k=\sum_{j=0}^k\binom{k}{j}(-1)^{k-j}\bE^{rj}$
combined with $\bE^{rj}\cdot\bE^{r-1}=\bE^{r(j+1)-1}$ gives the second
equality.
\end{proof}

\section{Power-sum expansion of $F_m$}

\begin{proposition}\label{prop:Fm}
For all $n\ge 0$ and $m\in\ZZ_{\ge 0}$,
\begin{equation}\label{eq:Fm}
   F_m \;=\; \sum_{\mu\vdash n}\;\frac{\eps(\mu)\,m^{\ell(\mu)}}{z_\mu}\;p_\mu.
\end{equation}
Consequently $m\mapsto F_m$ is a polynomial in $m$ of degree at most $n$
with values in $\Lam_n$, and its coefficient of $m^\ell$ is $P_\ell$.
\end{proposition}

\begin{proof}
Newton's identity in the plethystic form gives
$\log\bigl(\prod_i(1+x_i)\bigr)=\sum_{k\ge 1}\frac{(-1)^{k-1}}{k}\,p_k$,
so
\[
   \bE^m
   \;=\;\prod_i(1+x_i)^m
   \;=\;\exp\!\Bigl(m\sum_{k\ge 1}\tfrac{(-1)^{k-1}}{k}p_k\Bigr).
\]
Expanding the exponential and grouping by the resulting partition
$\mu=(1^{\ell_1}\,2^{\ell_2}\cdots)$,
\[
   \bE^m
   \;=\;\sum_{\ell_1,\ell_2,\ldots\ge 0}\;
        \prod_{j\ge 1}\frac{m^{\ell_j}(-1)^{(j-1)\ell_j}}{j^{\ell_j}\ell_j!}\,p_j^{\ell_j}
   \;=\;\sum_{\mu}\frac{m^{\ell(\mu)}\,\eps(\mu)}{z_\mu}\,p_\mu,
\]
using $\sum_j\ell_j(\mu)=\ell(\mu)$, $\sum_j(j-1)\ell_j(\mu)=n-\ell(\mu)$
(so the sign is $\eps(\mu)$), and $\prod_jj^{\ell_j}\ell_j!=z_\mu$.
Extracting the degree-$n$ part yields~\eqref{eq:Fm}. The final claim is
immediate: $m^{\ell(\mu)}$ has degree $\ell(\mu)\le n$ in $m$, and grouping
$F_m$ by the exponent $\ell$ of $m$ gives $F_m=\sum_{\ell=1}^{n}m^\ell\,P_\ell$
(the term $\ell=0$ vanishes for $n\ge 1$).
\end{proof}

\section{The subspace $W_n$ and the sign-obstruction}

\begin{proposition}\label{prop:Wn}
For $n\ge 1$, let
\[
   W_n \;:=\; \spann_{\QQ}\bigl\{P_\ell:1\le\ell\le n\bigr\}
   \;=\; \spann_{\QQ}\bigl\{F_m:m\ge 1\bigr\}
   \;\subset\;\Lam_n.
\]
Then $\dim W_n=n$. Moreover, for every $r\ge 1$,
\[
   \spann_{\QQ}\bigl\{\Frob(M_{n,r}^{(k)}):0\le k\le n\bigr\}
   \;=\; W_n.
\]
In particular $W_n^{\mathrm{BR}}=W_n$.
\end{proposition}

\begin{proof}
The partitions of $n$ split by length as a disjoint union
$\{\mu\vdash n\}=\bigsqcup_{\ell=1}^n\{\mu:\ell(\mu)=\ell\}$; every length
$\ell\in\{1,\dots,n\}$ is realised (by e.g.\ $(n-\ell+1,1^{\ell-1})$), so
each $P_\ell$ is nonzero. Since the $P_\ell$'s are supported on disjoint
subsets of the $p$-basis they are linearly independent, giving $\dim
W_n=n$.

Proposition~\ref{prop:Fm} gives $F_m=\sum_{\ell=1}^n m^\ell P_\ell$, so
$F_m\in W_n$ for every $m$; conversely the Vandermonde matrix
$\bigl(m^\ell\bigr)_{\ell,m}$ has rank $n$ on any $n$ distinct values of
$m\ge 1$, so $\{P_\ell\}$ lies in the $\QQ$-span of $\{F_m\}$. Thus the
two spans coincide.

Finally, for fixed $r\ge 1$, the values $r-1,2r-1,\dots,(n+1)r-1$ are $n+1$
distinct nonnegative integers. Proposition~\ref{prop:gf} gives
$\Frob(M_{n,r}^{(k)})=\sum_{j=0}^k\binom{k}{j}(-1)^{k-j}F_{r(j+1)-1}$; the
change-of-basis matrix from $\bigl(F_{r-1},F_{2r-1},\ldots,F_{(k+1)r-1}\bigr)$
to $\bigl(\Frob(M_{n,r}^{(0)}),\ldots,\Frob(M_{n,r}^{(k)})\bigr)$ is the
lower-triangular matrix $\bigl((-1)^{k-j}\binom{k}{j}\bigr)_{0\le j\le k}$,
which is invertible. Letting $k$ run to $n$ shows the two spans coincide.
Hence
$\spann\{\Frob(M_{n,r}^{(k)})\}_{0\le k\le n}=\spann\{F_{r-1},F_{2r-1},\ldots,F_{(n+1)r-1}\}=W_n$
(the last equality because $n+1$ values of $m$ suffice to span the
polynomial-in-$m$ image, cf.\ the Vandermonde argument above).
\end{proof}

\begin{proof}[Proof of Theorem~\ref{thm:main}]
By Proposition~\ref{prop:Wn} it suffices to prove $p_{(e^{k'})}\notin W_n$
under the hypothesis $k'\ge 2$, $e\ge 2$, and to exhibit the functional
$L_\nu$ in~\eqref{eq:L-explicit}.

Suppose for contradiction $p_{(e^{k'})}=\sum_{\ell=1}^n\alpha_\ell P_\ell$
for some $\alpha_\ell\in\QQ$. Since $P_\ell$ is supported on partitions of
length exactly $\ell$ and $p_{(e^{k'})}$ is supported at the single
partition $(e^{k'})$ (of length $k'$), we must have $\alpha_\ell=0$ for
$\ell\ne k'$; and $\alpha_{k'}\cdot\sum_{\ell(\mu)=k'}\frac{\eps(\mu)}{z_\mu}p_\mu
   \;=\;p_{(e^{k'})}$.
Reading off the coefficient of $p_{(e^{k'})}$ forces
$\alpha_{k'}\cdot\frac{\eps((e^{k'}))}{z_{(e^{k'})}}=1$, so
$\alpha_{k'}\ne 0$. Reading off the coefficient of $p_\nu$ for any
$\nu\vdash n$ with $\ell(\nu)=k'$, $\nu\ne(e^{k'})$, forces
$\alpha_{k'}\cdot\frac{\eps(\nu)}{z_\nu}=0$, i.e.\ $\alpha_{k'}=0$
(since $\eps(\nu),z_\nu\ne 0$). This is the desired contradiction \emph{as
soon as such a $\nu$ exists}, which under our hypotheses is witnessed by
$\nu=(e{+}1,\underbrace{e,\ldots,e}_{k'-2\text{ copies}},e{-}1)$: this has
length $k'\ge 2$, all parts $\ge 1$ (using $e\ge 2$), is weakly decreasing
(using $e{+}1\ge e\ge e{-}1$), sums to $ek'=n$, and is distinct from
$(e^{k'})$ because its first part is $e{+}1\ne e$.

For the explicit functional, define $L_\nu$ by~\eqref{eq:L-explicit}. It
reads only coefficients at partitions of length $k'$, hence vanishes on
$P_\ell$ for $\ell\ne k'$; on $P_{k'}$ it evaluates to
\begin{align*}
   L_\nu(P_{k'})
   &\;=\;\frac{z_{(e^{k'})}}{\eps((e^{k'}))}\cdot\frac{\eps((e^{k'}))}{z_{(e^{k'})}}
        \;-\;\frac{z_\nu}{\eps(\nu)}\cdot\frac{\eps(\nu)}{z_\nu} \;=\; 1-1 \;=\; 0.
\end{align*}
Hence $L_\nu$ annihilates all of $W_n=W_n^{\mathrm{BR}}$. It evaluates
on the target as
$L_\nu(p_{(e^{k'})})=z_{(e^{k'})}/\eps((e^{k'}))\ne 0$.
This completes the proof of Theorem~\ref{thm:main}.
\end{proof}

\section{The strong form: sign-ratio detectors}

The proof above uses only one functional per required contradiction, but
the full annihilator $W_n^\perp$ has dimension $p(n)-n$ and admits a basis
of ``sign-ratio detectors.''

\begin{corollary}\label{cor:strong}
Fix $n\ge 1$. For each length $\ell$ that supports two or more partitions
of $n$, choose a distinguished partition $\nu_\ell^{\mathrm{piv}}$ of $n$ of
length $\ell$ (the ``pivot''), and for each other partition $\mu$ of $n$ of
length $\ell$ define
\[
   L_{\mu,\ell}(\phi)
   \;:=\;\frac{z_{\mu}}{\eps(\mu)}\,[p_\mu]\phi
        \;-\;\frac{z_{\nu_\ell^{\mathrm{piv}}}}{\eps(\nu_\ell^{\mathrm{piv}})}\,[p_{\nu_\ell^{\mathrm{piv}}}]\phi.
\]
These $\{L_{\mu,\ell}\}$ form a basis of $W_n^\perp$; there are exactly
$p(n)-n$ of them.
\end{corollary}

\begin{proof}
Each $L_{\mu,\ell}$ reads coefficients only at length-$\ell$ partitions,
so vanishes on $P_{\ell'}$ for $\ell'\ne\ell$; on $P_\ell$ it evaluates as
$(z_\mu/\eps(\mu))\cdot(\eps(\mu)/z_\mu)-(z_{\nu_\ell^{\mathrm{piv}}}/\eps(\nu_\ell^{\mathrm{piv}}))\cdot(\eps(\nu_\ell^{\mathrm{piv}})/z_{\nu_\ell^{\mathrm{piv}}})=0$.
Hence $L_{\mu,\ell}\in W_n^\perp$. Linear independence: the
$L_{\mu,\ell}$'s use distinct pairs of $p_\mu$-coefficient-reads, and one
verifies by inspection that no nontrivial combination sums to zero. The
count is $\sum_\ell(\pi_\ell(n)-1)=p(n)-n$, where $\pi_\ell(n)$ is the
number of partitions of $n$ of length $\ell$, which matches
$\dim\Lam_n-\dim W_n$.
\end{proof}

\begin{remark}
The empirical Schur-basis functionals $L_1,L_2,L_3\in\Lam_6^*$ found in
the $(n,r)=(6,3)$ probe (see the accompanying probe log at
\texttt{probes/2026-08-09-bhattacharya-rhoades-wreath-superspace/RESULT.md})
are $\QQ$-linear combinations of the length-2 sign-ratio detectors:
\begin{align*}
L_1 &= \tfrac{8}{3}[s_{(4,2)}] - 8[s_{(3,3)}] + [s_{(3,2,1)}],\\
L_2 &= -[s_{(5,1)}] + 5[s_{(4,2)}] - 10[s_{(3,3)}] + [s_{(3,1^3)}],\\
L_3 &= [s_{(5,1)}] - \tfrac{20}{3}[s_{(4,2)}] + 15[s_{(3,3)}] - 5[s_{(2^3)}] + [s_{(2,1^4)}].
\end{align*}
These pair to $-18,-18,+18$ against $p_3^2$ respectively, in agreement with
the sign of the coefficient of the pivot $p_{(3,3)}$ in the target.
Computationally we verify that each $L_i$ annihilates all six $P_\ell$'s
for $n=6$; see the verification script \texttt{verify\_probe\_L.py}.
\end{remark}

\section{Verification}

The proof rests on three verifiable claims:
\begin{enumerate}[label=(V\arabic*)]
\item Proposition~\ref{prop:gf}: $\Frob(M_{n,r}^{(k)})=\sum_{j=0}^k\binom{k}{j}(-1)^{k-j}F_{r(j+1)-1}$.
\item Proposition~\ref{prop:Wn}: $\rank\spann\{\Frob(M_{n,r}^{(k)})\}_{k=0}^n=n$
      for each $r\ge 1$; equal to $\rank\spann\{F_m\}=\rank\spann\{P_\ell\}=n$.
\item Theorem~\ref{thm:main}: $p_e^{k'}\notin W_n$ for $k'\ge 2$, witnessed
      by an explicit functional.
\end{enumerate}
\noindent
The scripts in \verb|proofs/2026-08-09-br-work/| confirm all three:
\begin{itemize}
\item \verb|verify_gf.py|: for $(n,r)=(6,3)$ and each $k=0,\dots,6$,
       computes $\Frob(M_{n,r}^{(k)})$ two ways (via BR-orbit enumeration
       and via the formula) and checks equality in the $p$-basis. All
       $7/7$ match.
\item \verb|verify_span.py|: for each $n\in\{4,5,6,8,9,12\}$, computes
       $\rank\spann\{F_m\}=\rank\spann\{P_\ell\}=n$ (all cases pass, with
       union-span also equal to $n$); and for each of the $9$ test cases
       $(n,e)\in\{(6,2),(6,3),(8,2),(8,4),(9,3),(12,2),(12,3),(12,4),(12,6)\}$
       verifies $p_e^{n/e}\notin W_n$ and produces the explicit
       $L_\nu$-functional.
\item \verb|verify_probe_L.py|: for $n=6$, confirms each of the empirical
       $L_1,L_2,L_3$ from the probe evaluates to $0$ on all six $P_\ell$'s,
       and to $\pm 18$ on $p_3^2$.
\end{itemize}

\section*{Aesthetic connection to the seed}

Bhattacharya--Rhoades's wreath superspace, restricted to $S_n$, is a
generating series of \emph{permutation representations} governed by a
compact plethystic-exponential structure ($\bE^r-1$ per block). The
resulting slice-Frobenius family lies in an $n$-dimensional subspace of
$\Lam_n$ organised by \emph{cycle length alone}. This is one dimension per
possible length, and it is orthogonal to what $q_e^{(n)}=p_{(e^{k'})}$
demands: a delta at a single length-$k'$ partition, in a length class that
supports many partitions.

The obstruction is neither a Kostka-cone (nonnegativity) issue nor a
multiplicity-2 issue --- it is the incompatibility between a
\emph{length-graded} single-parameter family (BR side) and a
\emph{cycle-type-graded} point (composite-$d$ side). This matches the
Levi-stability coordinating-theme observed across 2025--2026: the BR
module category is ``one length'' below the resolution needed. The right
level of abstraction for the obstruction is the (only) $n$-dimensional
family $\{P_\ell\}_\ell$, and the sign-ratio detectors capture the
codimension-$(p(n)-n)$ complement exactly.

\medskip

\emph{This completes the proof of the Bhattacharya--Rhoades sign-obstruction theorem.}

\begin{thebibliography}{9}
\bibitem{BR} A.\ Bhattacharya, B.\ Rhoades.
    \emph{The wreath superspace and the fermionic-degree filtration.}
    arXiv:2606.30977, 2026.
\bibitem{ClioSign} Clio.
    \emph{Virtual-character sign-pattern theorem for the composite-$d$ symmetric function $q_e^{(n)}=p_e^{k'}$.}
    Standalone proof, 2026-08-08.
\end{thebibliography}

\end{document}
